\begin{table}[t]
\centering
\small
\begin{tabular}{L{0.20\linewidth} L{0.34\linewidth} L{0.36\linewidth}}
\toprule
Dimension & Source paper (2018) & Upgraded paper (2026 framing) \\
\midrule
Primary question & Can 2D and 1D CNNs compete with 3D CNNs after Hilbert-based mapping? & When should serialization be treated as a first-class interface for irregular 3D and biomolecular learning? \\
Paper type & Empirical methods note & Hybrid interface paper with an executable probe \\
Evidence base & ModelNet10 and K-Ras / H-Ras experiments & Source-paper analysis, modern primary literature, and a deterministic context-budget study \\
Main contribution & A reversible 3D-to-2D / 1D mapping strategy & A taxonomy, executable locality study, design framework, and biomolecular research agenda for serialized structural learning \\
Acceptance story & Better efficiency than volumetric baselines under a 2018 benchmark setup & A timely interface paper that connects early Hilbert mappings to modern sequence-native 3D learning and demonstrates why ordering quality matters under bounded context \\
\bottomrule
\end{tabular}
\caption{Why the new manuscript is not a facelift. The central claim, paper type, evidence base, and contribution all change materially.}
\label{tab:source-vs-upgrade}
\end{table}
